% !TeX program = lualatex
% ================================================================
% Polish Parallel Text Version of SequencesNthTermIdeasStarters
%
% Generated by the server using the following settings:
%
%   language            Polish (pol)
%   text direction      left to right
%   font                Noto Serif
%   fallback font       Noto Serif
%   built from          latex/starters/targeted/KS3/sequences/sequencesNthTermIdeasStarters.tex
%   vocabulary key      latex/starters/targeted/KS3/sequences/sequencesNthTermIdeasStarters/L2/pol/sequencesNthTermIdeasStarters_polishKey.tex
%   tier 3 vocabulary   green   RGB 0 128 0
%   English text        black   RGB 0 0 0
%   Polish text         purple  RGB 112 48 160
%   layout macros       version 3
%   sheet layout        version 1
%
% Please don't edit any information in this comment block - the server
% compares it against the current settings and rebuilds this file when
% they differ, so changes here would be overwritten. However, feel free
% to make updates to the LaTeX below if you want to edit it.
% ================================================================

\documentclass{beamer}
% ================================================================
% Copied from the English sheet this was translated from, so that
% anything it sets up for itself still works here
% ================================================================
\usepackage{tikz}
\usepackage{amsmath}

% ---- Minimal styling ----
\setbeamertemplate{navigation symbols}{}
\setbeamertemplate{footline}{}
\setbeamercolor{background canvas}{bg=white}
\setbeamercolor{frametitle}{fg=black}
\setbeamerfont{frametitle}{size=\Large, series=\bfseries}
\usefonttheme{professionalfonts}

% ---- Global tikz baseline ----
\tikzset{every picture/.style={baseline=(current bounding box.south)}}

% ================================================================
% TikZ pattern commands
% ================================================================

% Matchstick squares: n squares in a row
\newcommand{\matchsquares}[2][black]{%
  \begin{tikzpicture}[scale=0.60]
    \pgfmathsetmacro\n{#2}
    \color{#1}
    \foreach \x in {0,...,\n} {
      \draw[very thick, cap=round] (\x,0) -- (\x,1);
    }
    \foreach \x in {0,...,\the\numexpr\n-1\relax} {
      \draw[very thick, cap=round] (\x,0)   -- (\x+1,0);
      \draw[very thick, cap=round] (\x,1)   -- (\x+1,1);
    }
  \end{tikzpicture}%
}

% Matchstick triangles: chain of n triangles
\newcommand{\matchtriangles}[2][black]{%
  \begin{tikzpicture}[scale=0.60]
    \pgfmathsetmacro\n{#2}
    \pgfmathsetmacro\h{0.866}
    \color{#1}
    \draw[very thick] (0,0) -- (\n,0);
    \draw[very thick] (0,0)
      \foreach \k in {1,...,\n} {
        -- ({\k - 0.5}, {\h}) -- (\k,0)
      };
  \end{tikzpicture}%
}

% Rectangular dots: n columns x (n+1) rows
\newcommand{\rectdots}[2][black]{%
  \begin{tikzpicture}[scale=0.60]
    \pgfmathsetmacro\n{#2}
    \color{#1}
    \foreach \i in {0,...,\the\numexpr\n-1\relax} {
      \foreach \j in {0,...,\n} {
        \fill (0.4*\i, 0.4*\j) circle (0.07);
      }
    }
  \end{tikzpicture}%
}

% Cross / plus dots
\newcommand{\crossdots}[2][black]{%
  \begin{tikzpicture}[scale=0.60]
    \pgfmathsetmacro\n{#2}
    \color{#1}
    \fill (0,0) circle (0.07);
    \ifnum\n>1
      \foreach \d in {1,...,\the\numexpr\n-1\relax} {
        \fill (0,   0.4*\d) circle (0.07);
        \fill (0,  -0.4*\d) circle (0.07);
        \fill ( 0.4*\d, 0)  circle (0.07);
        \fill (-0.4*\d, 0)  circle (0.07);
      }
    \fi
  \end{tikzpicture}%
}

% 4-times-table dots: n rows of 4 dots
\newcommand{\fourtabledots}[2][black]{%
  \begin{tikzpicture}[scale=0.60]
    \pgfmathsetmacro\n{#2}
    \color{#1}
    \foreach \row in {0,...,\the\numexpr\n-1\relax} {
      \foreach \col in {0,1,2,3} {
        \fill (0.4*\col, 0.4*\row) circle (0.07);
      }
    }
  \end{tikzpicture}%
}

% ================================================================
% Answer-overlay helpers
% ================================================================
\newcommand{\ablank}[1]{%
  \alt<2>{\textcolor{red}{#1}}{\underline{\phantom{#1}}}%
}
\newcommand{\ashow}[1]{\uncover<2->{\textcolor{red}{\small #1}}}
\newcommand{\ashowq}[1]{\alt<2->{\textcolor{red}{\small #1}}{?}}

% Answers drawn straight on to a TikZ diagram, revealed on overlay 2 like \ashow.
% Space is reserved on overlay 1, so the diagram does not move between slides.
\usetikzlibrary{overlay-beamer-styles}
\tikzset{
  ans/.style     = {red, thick, visible on=<2->},
  ansfill/.style = {red!15, visible on=<2->},
  anslab/.style  = {red, font=\tiny, visible on=<2->},
}

% ================================================================

% Characters this language's font has not got are borrowed from another,
% rather than being silently dropped - see docs/translated-sheet-latex.md
\directlua{%
  if luaotfload and luaotfload.add_fallback then
    luaotfload.add_fallback("eallatinfallback",
      { "Noto Serif:mode=harf;" })
  else
    tex.error("this luaotfload is too old to register a fallback font")
  end
}
\usepackage[english]{babel}
\babelprovide[import]{polish}
\babelfont[polish]{rm}[Renderer=Harfbuzz, BoldFont={Noto Serif}, ItalicFont={Noto Serif}, BoldItalicFont={Noto Serif}, RawFeature={fallback=eallatinfallback}]{Noto Serif}
\babelfont[polish]{sf}[Renderer=Harfbuzz, BoldFont={Noto Serif}, ItalicFont={Noto Serif}, BoldItalicFont={Noto Serif}, RawFeature={fallback=eallatinfallback}]{Noto Serif}
\newcommand{\ealtext}[1]{\foreignlanguage{polish}{#1}}
\newcommand{\ealtextblock}[1]{{\raggedright\foreignlanguage{polish}{#1}\par}}
% ================================================================
% EAL layout helpers
% ================================================================
\usepackage{xcolor}

\definecolor{ealkeycolour}{RGB}{0,128,0}
\definecolor{ealenglish}{RGB}{0,0,0}
\definecolor{ealtwocolour}{RGB}{112,48,160}

\newsavebox{\ealboxen}
\newsavebox{\ealboxtr}
\newlength{\ealglosswd}
\newlength{\ealglossraise}
\setlength{\ealglossraise}{1.15em}

% a tier 3 word, in English and in translation alike
\newcommand{\ealkey}[1]{\textcolor{ealkeycolour}{#1}}
\newcommand{\ealkeytr}[1]{\textcolor{ealkeycolour}{\ealtext{#1}}}

% One translated word above one English word. Built from kernel
% primitives, never a tabular, which would break wherever the sheet
% loads array, booktabs or tabularx. The box is as wide as the wider of
% the two, so a long translation pushes its neighbours aside instead of
% overlapping them, and it claims its full height so a table row or a
% TikZ node makes room for it.
\newcommand{\ealgloss}[2]{%
  \sbox{\ealboxen}{\ealkey{#1}}%
  \sbox{\ealboxtr}{\small\textcolor{ealtwocolour}{\ealtext{#2}}}%
  \setlength{\ealglosswd}{\wd\ealboxen}%
  \ifdim\wd\ealboxtr>\ealglosswd\setlength{\ealglosswd}{\wd\ealboxtr}\fi
  \leavevmode
  \makebox[\ealglosswd][c]{%
    \makebox[0pt][c]{\usebox{\ealboxen}}%
    \makebox[0pt][c]{\raisebox{\ealglossraise}{\usebox{\ealboxtr}}}%
  }%
}

% opens up line spacing around a block of glosses, so the raised
% translations sit clear of the line above
\newenvironment{ealglossed}{\par\linespread{2.1}\selectfont\ignorespaces}{\par}

% a whole translated sentence above its English counterpart. block level,
% so it does not belong inside a tabular cell
\newcommand{\ealpara}[2]{%
  \par\addvspace{0.5\baselineskip}%
  {\color{ealtwocolour}\ealtextblock{#1}}%
  \penalty200
  {\color{ealenglish}#2\par}%
  \addvspace{0.5\baselineskip}%
}

\begin{document}
% ================================================================

% ---- Title ----
\begin{frame}[plain]
  \centering
  {\LARGE\bfseries \ealkey{Sequences} Starters}\\[0.5em]
\end{frame}

% ================================================================
% STARTER 1
% ================================================================
\begin{frame}[shrink=5]{Starter 1 (1 of 3)}
  \begin{minipage}[t]{0.48\textwidth}
    \small
    \ealpara{Q1. Wzór z zapałek}{\textbf{Q1.\enspace Matchstick Pattern}}
    \matchsquares[red]{1}\quad
    \matchsquares[blue]{2}\quad
    \matchsquares[green!70!black]{3}\quad
    \ashowq{\matchsquares[orange]{4}}
    \ealpara{(a) Narysuj 4. kształt.}{(a)\;Draw the 4th shape.}
    \ealpara{(b) Ile zapałek używa?}{(b)\;How many matchsticks does it use?}
    \ashow{\ealpara{(b) 13 zapałek}{(b)\;13 matchsticks}}
  \end{minipage}\hfill
  \begin{minipage}[t]{0.48\textwidth}
    \small
    \ealpara{Q2. Uzupełnij \ealkeytr{ciąg}}{\textbf{Q2.\enspace Complete the \ealkey{Sequence}}}
    \footnotesize
    5,\;10,\;15,\;20,\;
    \ablank{25},\;\ablank{30},\;\ablank{35}
    \ashow{\ealpara{(dodaj 5 za każdym razem)}{(add 5 each time)}}
  \end{minipage}
\end{frame}

\begin{frame}[shrink=5]{Starter 1 (2 of 3)}
  \small
  \ealpara{Q3. Zapisz pierwsze 5 \ealkeytr{wyrazów} \ealkeytr{ciągu} o \ealkeytr{wzorze na n-ty wyraz ciągu}: $2n$}{\textbf{Q3.\enspace Write the first 5 \ealkey{terms} of the \ealkey{sequence} with \ealkey{position-to-term rule}: $\boldsymbol{2n}$}}
  \footnotesize
  \ablank{2},\;\ablank{4},\;\ablank{6},\;\ablank{8},\;\ablank{10}

  \vspace{0.45cm}

  \small
  \ealpara{Sam zapisuje tabliczkę mnożenia przez 3: $3, 6, 9, 12, 15\ldots$ Następnie \textbf{dodaje 2} do każdego \ealkeytr{wyrazu}. Zapisz pierwsze 5 \ealkeytr{wyrazów} \ealkeytr{ciągu} Sama. Jaki jest \ealkeytr{wzór na n-ty wyraz}?}{\textbf{Q4.\enspace} Sam writes the 3 \ealkey{times table}: $3, 6, 9, 12, 15\ldots$ He then \textbf{adds 2} to every \ealkey{term}. Write the first 5 \ealkey{terms} of Sam's \ealkey{sequence}. What is the \ealkey{$n$th term} rule?}
  \ashow{\ealpara{Wyrazy: $5, 8, 11, 14, 17$. \quad \ealkeytr{wzór na n-ty wyraz}: $3n+2$}{\ealkey{Terms}: $5, 8, 11, 14, 17$. \quad \ealkey{$n$th term}: $\boldsymbol{3n+2}$}}
\end{frame}

\begin{frame}[shrink=5]{Starter 1 (3 of 3)}
  \begin{minipage}[t]{0.48\textwidth}
    \small
    \ealpara{Rozszerzenie Q1}{\textbf{Q5.\enspace Extending Q1}}
    \matchsquares[red]{1}\quad
    \matchsquares[blue]{2}\quad
    \matchsquares[green!70!black]{3}\quad
    \ashowq{\matchsquares[orange]{4}}\\[0.4em]
    \footnotesize
    \ealpara{Używając wzoru z zapałek z Q1:}{Using the matchstick pattern from Q1:}
    \ealpara{(a) Ile zapałek w 10. kształcie?}{(a)\;How many matchsticks in the \textbf{10th} shape?}
    \ealpara{(b) Zapisz regułę dla n-tego kształtu.}{(b)\;Write a rule for the $n$th shape.}
    \ashow{\ealpara{(a) 31 zapałek.}{(a)\;31 matchsticks.}}
    \ashow{\ealpara{(b) \ealkeytr{wzór na n-ty wyraz ciągu}: $3n+1$}{(b)\;\ealkey{position-to-term rule}:\;$3n+1$}}
  \end{minipage}\hfill
  \begin{minipage}[t]{0.48\textwidth}
    \small
    \ealpara{Q6. Wyzwanie}{\textbf{Q6.\;Challenge}}
    \footnotesize
    \ealpara{Czy $\boldsymbol{100}$ jest \ealkeytr{wyrazem} \ealkeytr{ciągu} o \ealkeytr{wzorze na n-ty wyraz} $3n+1$?}{Is $\boldsymbol{100}$ a \ealkey{term} in the \ealkey{sequence} with \ealkey{$n$th term} $3n+1$\,?}
    \ealpara{Pokaż, skąd to wiesz.}{Show how you know.}
    \ashow{\ealpara{Tak, gdy $n=33$,}{Yes, when $n=33$,}}
    \ashow{$3n+1=3\times33+1=99+1=100$}
  \end{minipage}
\end{frame}

% ================================================================
% STARTER 2
% ================================================================
\begin{frame}[shrink=5]{Starter 2 (1 of 3)}
  \begin{minipage}[t]{0.48\textwidth}
    \small
    \ealpara{Q1. Wzór kropek}{\textbf{Q1.\enspace Dot Pattern}}
    \fourtabledots[red]{1}\quad
    \fourtabledots[blue]{2}\quad
    \fourtabledots[green!70!black]{3}\quad
    \ashowq{\fourtabledots[orange]{4}}\\[0.4em]
    \footnotesize
    \ealpara{(a) Narysuj 4. obrazek.}{(a)\;Draw the 4th picture.}
    \ealpara{(b) Ile kropek zawiera?}{(b)\;How many dots does it contain?}
    \ashow{\ealpara{(b) 16 kropek}{(b)\;16 dots}}
  \end{minipage}\hfill
  \begin{minipage}[t]{0.48\textwidth}
    \small
    \ealpara{Q2. Uzupełnij \ealkeytr{ciąg}}{\textbf{Q2.\enspace Complete the \ealkey{Sequence}}}
    \footnotesize
    4,\;8,\;12,\;16,\;
    \ablank{20},\;\ablank{24},\;\ablank{28}
    \ashow{\ealpara{(dodaj 4 za każdym razem)}{(add 4 each time)}}
  \end{minipage}
\end{frame}

\begin{frame}[shrink=5]{Starter 2 (2 of 3)}
  \small
  \ealpara{Q3. Zapisz pierwsze 5 \ealkeytr{wyrazów} $3n+1$}{\textbf{Q3.\enspace Write the first 5 \ealkey{terms} of $\boldsymbol{3n+1}$}}
  \footnotesize
  \ablank{4},\;\ablank{7},\;\ablank{10},\;\ablank{13},\;\ablank{16}

  \vspace{0.45cm}

  \small
  \ealpara{Alex zapisuje tabliczkę mnożenia przez 5: $5, 10, 15, 20, 25\ldots$ Następnie \textbf{odejmuje 1} od każdego \ealkeytr{wyrazu}. Zapisz pierwsze 5 \ealkeytr{wyrazów} \ealkeytr{ciągu} Alexa. Jaki jest \ealkeytr{wzór na n-ty wyraz}?}{\textbf{Q4.\enspace} Alex writes the 5 \ealkey{times table}: $5, 10, 15, 20, 25\ldots$ He then \textbf{subtracts 1} from every \ealkey{term}. Write the first 5 \ealkey{terms} of Alex's \ealkey{sequence}. What is the \ealkey{$n$th term} rule?}
  \ashow{\ealpara{Wyrazy: $4, 9, 14, 19, 24$. \quad \ealkeytr{wzór na n-ty wyraz}: $5n-1$}{\ealkey{Terms}: $4, 9, 14, 19, 24$. \quad \ealkey{$n$th term}: $\boldsymbol{5n-1}$}}
\end{frame}

\begin{frame}[shrink=5]{Starter 2 (3 of 3)}
  \begin{minipage}[t]{0.48\textwidth}
    \small
    \ealpara{Rozszerzenie Q1}{\textbf{Q5.\enspace Extend Q1}}
    \fourtabledots[red]{1}\quad
    \fourtabledots[blue]{2}\quad
    \fourtabledots[green!70!black]{3}\quad
    \ashowq{\fourtabledots[orange]{4}}\\[0.4em]
    \footnotesize
    \ealpara{Używając swojego wzoru kropek z Q1:}{Using your dot pattern from Q1:}
    \ealpara{(a) Ile kropek w 10. obrazku?}{(a)\;How many dots in the \textbf{10th} picture?}
    \ealpara{(b) Zapisz regułę dla n-tego obrazka.}{(b)\;Write a rule for the $n$th picture.}
    \ashow{\ealpara{(a) 40 kropek. \quad (b) Reguła: $4n$}{(a)\;40 dots.\quad (b)\;Rule:\;$4n$}}
  \end{minipage}\hfill
  \begin{minipage}[t]{0.48\textwidth}
    \small
    \ealpara{Q6. Wyzwanie}{\textbf{Q6.\;Challenge}}
    \footnotesize
    \ealpara{\ealkeytr{Ciąg} A ma \ealkeytr{wzór na n-ty wyraz} $4n-2$.}{\ealkey{Sequence} A has \ealkey{$n$th term} $4n-2$.}
    \ealpara{\ealkeytr{Ciąg} B ma \ealkeytr{wzór na n-ty wyraz} $2n+6$.}{\ealkey{Sequence} B has \ealkey{$n$th term} $2n+6$.}
    \ealpara{Czy mają wspólny \ealkeytr{wyraz} na tej samej pozycji?}{Do they share a \ealkey{term} \emph{at the same position}?}
    \ashow{\ealpara{$4n-2=2n+6\;\Rightarrow\;n=4.$ Oba równe \textbf{14} na pozycji $n=4$.}{$4n-2=2n+6\;\Rightarrow\;n=4.$\; Both equal\;\textbf{14} at position $n=4$.}}
  \end{minipage}
\end{frame}

% ================================================================
% STARTER 3
% ================================================================
\begin{frame}[shrink=5]{Starter 3 (1 of 3)}
  \begin{minipage}[t]{0.48\textwidth}
    \small
    \ealpara{Q1. Trójkąty z zapałek}{\textbf{Q1.\enspace Matchstick Triangles}}
    \matchtriangles[red]{1}\quad
    \matchtriangles[blue]{2}\quad
    \matchtriangles[green!70!black]{3}\quad
    \ashowq{\matchtriangles[orange]{4}}\\[0.4em]
    \footnotesize
    \ealpara{(a) Narysuj 4. kształt.}{(a)\;Draw the 4th shape.}
    \ealpara{(b) Ile zapałek używa?}{(b)\;How many matchsticks does it use?}
    \ashow{\ealpara{(b) 12 zapałek}{(b)\;12 matchsticks}}
  \end{minipage}\hfill
  \begin{minipage}[t]{0.48\textwidth}
    \small
    \ealpara{Q2. Uzupełnij \ealkeytr{ciąg}}{\textbf{Q2.\enspace Complete the \ealkey{Sequence}}}
    \footnotesize
    7,\;14,\;21,\;28,\;
    \ablank{35},\;\ablank{42},\;\ablank{49}
    \ashow{\ealpara{(dodaj 7 za każdym razem --- 7$\times$ \ealkeytr{tabliczka mnożenia})}{(add 7 each time --- the 7$\times$ \ealkey{times table})}}
  \end{minipage}
\end{frame}

\begin{frame}[shrink=5]{Starter 3 (2 of 3)}
  \small
  \ealpara{Q3. Zapisz pierwsze 5 \ealkeytr{wyrazów} $5n-2$}{\textbf{Q3.\enspace Write the first 5 \ealkey{terms} of $\boldsymbol{5n-2}$}}
  \footnotesize
  \ablank{3},\;\ablank{8},\;\ablank{13},\;\ablank{18},\;\ablank{23}

  \vspace{0.45cm}

  \small
  \ealpara{\ealkeytr{Ciąg} \textbf{zaczyna się od 5} i \textbf{rośnie o 3} za każdym razem. Zapisz pierwsze 5 \ealkeytr{wyrazów}. Jaki jest \ealkeytr{wzór na n-ty wyraz}?}{\textbf{Q4.\enspace} A \ealkey{sequence} \textbf{starts at 5} and \textbf{increases by 3} each time. Write the first 5 \ealkey{terms}. What is the \ealkey{$n$th term} rule?}
  \ashow{\ealpara{Wyrazy: $5, 8, 11, 14, 17$. \quad \ealkeytr{wzór na n-ty wyraz}: $3n+2$}{\ealkey{Terms}: $5, 8, 11, 14, 17$. \quad \ealkey{$n$th term}: $\boldsymbol{3n+2}$}}
\end{frame}

\begin{frame}[shrink=5]{Starter 3 (3 of 3)}
  \begin{minipage}[t]{0.48\textwidth}
    \small
    \ealpara{Rozszerzenie Q1}{\textbf{Q5.\enspace Extending Q1}}
    \matchtriangles[red]{1}\quad
    \matchtriangles[blue]{2}\quad
    \matchtriangles[green!70!black]{3}\quad
    \ashowq{\matchtriangles[orange]{4}}\\[0.4em]
    \footnotesize
    \ealpara{Używając swojego wzoru trójkątów z Q1:}{Using your triangle pattern from Q1:}
    \ealpara{(a) Zapisz regułę dla n-tego kształtu.}{(a)\;Write a rule for the $n$th shape.}
    \ealpara{(b) Ile zapałek w 20. kształcie?}{(b)\;How many matchsticks in the \textbf{20th} shape?}
    \ashow{\ealpara{(a) Reguła: $3n$. \quad (b) 60 zapałek.}{(a)\;Rule:\;$3n$.\quad (b)\;60 matchsticks.}}
  \end{minipage}\hfill
  \begin{minipage}[t]{0.48\textwidth}
    \small
    \ealpara{Q6. Wyzwanie}{\textbf{Q6.\;Challenge}}
    \footnotesize
    \ealpara{\ealkeytr{wzór na n-ty wyraz} \ealkeytr{ciągu} to $4n+3$.}{\ealkey{The $n$th term} of a \ealkey{sequence} is $4n+3$.}
    \ealpara{Który \ealkeytr{wyraz} jest \textbf{pierwszym, który przekracza 75}?}{Which \ealkey{term} is the \textbf{first to exceed 75}?}
    \ealpara{Pokaż wyraźnie swoje obliczenia.}{Show your working clearly.}
    \ashow{\ealpara{$4n+3>75\;\Rightarrow\;n>18.$ 19. \ealkeytr{wyraz}: $4(19)+3=\boldsymbol{79}$.}{$4n+3>75\;\Rightarrow\;n>18.$\; 19th \ealkey{term}: $4(19)+3=\boldsymbol{79}.$}}
  \end{minipage}
\end{frame}

% ================================================================
% STARTER 4
% ================================================================
\begin{frame}[shrink=5]{Starter 4 (1 of 3)}
  \begin{minipage}[t]{0.48\textwidth}
    \small
    \ealpara{Q1. Kropki}{\textbf{Q1.\enspace Dots}}
    \rectdots[red]{1}\quad
    \rectdots[blue]{2}\quad
    \rectdots[green!70!black]{3}\quad
    \ashowq{\rectdots[orange]{4}}\\[0.4em]
    \footnotesize
    \ealpara{(a) Narysuj 4. wzór.}{(a)\;Draw the 4th pattern.}
    \ealpara{(b) Ile kropek zawiera?}{(b)\;How many dots does it contain?}
    \ashow{\ealpara{(b) 20 kropek}{(b)\;20 dots}}
  \end{minipage}\hfill
  \begin{minipage}[t]{0.48\textwidth}
    \small
    \ealpara{Q2. Uzupełnij \ealkeytr{ciąg}}{\textbf{Q2.\enspace Complete the \ealkey{Sequence}}}
    \footnotesize
    1,\;5,\;9,\;13,\;
    \ablank{17},\;\ablank{21},\;\ablank{25}
    \ashow{\ealpara{(dodaj 4 za każdym razem)}{(add 4 each time)}}
  \end{minipage}
\end{frame}

\begin{frame}[shrink=5]{Starter 4 (2 of 3)}
  \small
  \ealpara{Q3. Zapisz pierwsze 5 \ealkeytr{wyrazów} $10-2n$}{\textbf{Q3.\enspace Write the first 5 \ealkey{terms} of $\boldsymbol{10-2n}$}}
  \footnotesize
  \ablank{8},\;\ablank{6},\;\ablank{4},\;\ablank{2},\;\ablank{0}

  \vspace{0.45cm}

  \small
  \ealpara{Mia zapisuje tabliczkę mnożenia przez 6: $6, 12, 18, 24, 30\ldots$ Następnie \textbf{odejmuje 4} od każdego \ealkeytr{wyrazu}. Zapisz pierwsze 5 \ealkeytr{wyrazów}, znajdź \ealkeytr{wzór na n-ty wyraz} i znajdź 50. \ealkeytr{wyraz}.}{\textbf{Q4.\enspace} Mia writes the 6 \ealkey{times table}: $6, 12, 18, 24, 30\ldots$ She then \textbf{subtracts 4} from every \ealkey{term}. Write the first 5 \ealkey{terms}, find the \ealkey{$n$th term}, and find the 50th \ealkey{term}.}
  \ashow{\ealpara{Wyrazy: $2, 8, 14, 20, 26$. \quad \ealkeytr{wzór na n-ty wyraz}: $6n-4$. \quad 50. \ealkeytr{wyraz}: $296$.}{\ealkey{Terms}: $2, 8, 14, 20, 26$.\quad \ealkey{$n$th term}: $6n-4$.\quad 50th \ealkey{term}: $\boldsymbol{296}$.}}
\end{frame}

\begin{frame}[shrink=5]{Starter 4 (3 of 3)}
  \begin{minipage}[t]{0.48\textwidth}
    \small
    \ealpara{Q5. Znajdź \ealkeytr{wzór na n-ty wyraz}}{\textbf{Q5.\enspace Find the \ealkey{$n$th Term}}}
    \footnotesize
    $1,\;5,\;9,\;13,\;17,\;\ldots$
    \ealpara{Jaki jest \ealkeytr{wzór na n-ty wyraz}?}{What is the \ealkey{$n$th term} rule?}
    \ashow{\ealpara{\ealkeytr{Różnice}: $+4$. \quad Pierwszy \ealkeytr{wyraz}: 1. \quad Reguła: $4n-3$}{\ealkey{Differences}: $+4$. \;First \ealkey{term}: 1. \quad Rule:\;$\boldsymbol{4n-3}$}}
  \end{minipage}\hfill
  \begin{minipage}[t]{0.48\textwidth}
    \small
    \ealpara{Q6. Wyzwanie}{\textbf{Q6.\;Challenge}}
    \footnotesize
    \ealpara{\ealkeytr{wzór na n-ty wyraz} \ealkeytr{ciągu} to $10-2n$.}{\ealkey{The $n$th term} of a \ealkey{sequence} is $10-2n$.}
    \ealpara{Który to \textbf{ostatni} \ealkeytr{wyraz} będący \ealkeytr{liczbą dodatnią}?}{Which is the \textbf{last \ealkey{positive} \ealkey{term}}?}
    \ealpara{Jaki jest \textbf{pierwszy} \ealkeytr{ujemny} \ealkeytr{wyraz}?}{What is the \textbf{first \ealkey{negative} \ealkey{term}}?}
    \ashow{\ealpara{$n=4$: $+2$ (ostatnia \ealkeytr{liczba dodatnia}).}{$n=4$:\;$+2$ (last \ealkey{positive}).}}
    \ashow{\ealpara{$n=5$: $0$.}{$n=5$:\;$0$.}}
    \ashow{\ealpara{$n=6$: $-2$ (pierwszy \ealkeytr{ujemny}).}{$n=6$:\;$-2$ (first \ealkey{negative}).}}
  \end{minipage}
\end{frame}

% ================================================================
% STARTER 5
% ================================================================
\begin{frame}[shrink=5]{Starter 5 (1 of 3)}
  \begin{minipage}[t]{0.48\textwidth}
    \small
    \ealpara{Q1. Kropki}{\textbf{Q1.\enspace Dots}}
    \crossdots[red]{1}\quad
    \crossdots[blue]{2}\quad
    \crossdots[green!70!black]{3}\quad
    \ashowq{\crossdots[orange]{4}}\\[0.4em]
    \footnotesize
    \ealpara{(a) Narysuj 4. kształt.}{(a)\;Draw the 4th shape.}
    \ealpara{(b) Ile kropek w kształcie 5?}{(b)\;How many dots in shape 5?}
    \ashow{\ealpara{(b) Kształt 5 ma 17 kropek}{(b)\;Shape 5 has 17 dots}}
  \end{minipage}\hfill
  \begin{minipage}[t]{0.48\textwidth}
    \small
    \ealpara{Q2. Uzupełnij \ealkeytr{ciąg}}{\textbf{Q2.\enspace Complete the \ealkey{Sequence}}}
    \footnotesize
    1,\;4,\;9,\;16,\;
    \ablank{25},\;\ablank{36},\;\ablank{49}
    \ashow{\ealpara{(\ealkeytr{liczby kwadratowe}, $n^2$)}{(\ealkey{square numbers},\;$n^2$)}}
  \end{minipage}
\end{frame}

\begin{frame}[shrink=5]{Starter 5 (2 of 3)}
  \small
  \ealpara{Q3. Zapisz pierwsze 5 \ealkeytr{wyrazów} $2n+5$}{\textbf{Q3.\enspace Write the first 5 \ealkey{terms} of $\boldsymbol{2n+5}$}}
  \footnotesize
  \ablank{7},\;\ablank{9},\;\ablank{11},\;\ablank{13},\;\ablank{15}

  \vspace{0.45cm}

  \small
  \ealpara{Q4. Znajdź \ealkeytr{wzór na n-ty wyraz}}{\textbf{Q4.\enspace Find the \ealkey{$n$th Term}}}
  \footnotesize
  (a)\;$3,\;5,\;7,\;9,\;11,\;\ldots$\quad
  \ashow{$2n+1$}\\[0.3em]
  (b)\;$6,\;10,\;14,\;18,\;22,\;\ldots$\quad
  \ashow{$4n+2$}
\end{frame}

\begin{frame}[shrink=5]{Starter 5 (3 of 3)}
  \begin{minipage}[t]{0.48\textwidth}
    \small
    \ealpara{Q5.}{\textbf{Q5.\enspace}}
    \footnotesize
    \ealpara{Jake mówi: ``Zaczynam od 3 i dodaję 4 za każdym razem.''}{Jake says: ``I start at 3 and add 4 each time.''}
    \ealpara{(a) Zapisz jego pierwsze 5 \ealkeytr{wyrazów}, (b) znajdź \ealkeytr{wzór na n-ty wyraz}, i (c) zdecyduj, czy $\boldsymbol{59}$ jest w jego \ealkeytr{ciągu}. Pokaż swoje rozumowanie.}{(a) Write his first 5 \ealkey{terms}, (b) find the \ealkey{$n$th term} rule, and (c) decide whether $\boldsymbol{59}$ is in his \ealkey{sequence}. Show your reasoning.}
    \ashow{\ealpara{(a) $3,7,11,15,19$.}{(a) $3,7,11,15,19$.}}
    \ashow{\ealpara{(b) \ealkeytr{wzór na n-ty wyraz}: $4n-1$.}{(b) \ealkey{$n$th term}: $4n-1$.}}
    \ashow{\ealpara{(c) $4n-1=59\;\Rightarrow\;n=15.$ \textbf{Tak!} (15. \ealkeytr{wyraz})}{(c) $4n-1=59\;\Rightarrow\;n=15.$\;\textbf{Yes!} (15th \ealkey{term})}}
  \end{minipage}\hfill
  \begin{minipage}[t]{0.48\textwidth}
    \small
    \ealpara{Q6. Wyzwanie}{\textbf{Q6.\;Challenge}}
    \footnotesize
    \ealpara{\ealkeytr{Ciąg} Mai ma \ealkeytr{wzór na n-ty wyraz} $3n+1$.}{\ealkey{Mia's sequence} has \ealkey{$n$th term} $3n+1$.}
    \ealpara{\ealkeytr{Ciąg} Jake'a ma \ealkeytr{wzór na n-ty wyraz} $2n+6$.}{\ealkey{Jake's sequence} has \ealkey{$n$th term} $2n+6$.}
    \ealpara{Znajdź \textbf{pierwsze trzy wartości}, które pojawiają się w \emph{obu} \ealkeytr{ciągach}. Co zauważasz?}{Find the \textbf{first three values} that appear in \emph{both} \ealkey{sequences}. What do you notice about them?}
    \ashow{\ealpara{Wspólne wartości: $10, 16, 22, \ldots$ (rosną o 6 za każdym razem).}{Shared values: $10, 16, 22, \ldots$ (they increase by 6 each time).}}
  \end{minipage}
\end{frame}

% ================================================================
% STARTER 6
% ================================================================
\begin{frame}[shrink=5]{Starter 6 (1 of 3)}
  \begin{minipage}[t]{0.48\textwidth}
    \small
    \ealpara{Q1. Trójkąty z zapałek}{\textbf{Q1.\enspace Matchstick Triangles}}
    \matchtriangles[red]{1}\quad
    \matchtriangles[blue]{2}\quad
    \matchtriangles[green!70!black]{3}\quad
    \ashowq{\matchtriangles[orange]{4}}\\[0.4em]
    \footnotesize
    \ealpara{(a) Narysuj 4. kształt.}{(a)\;Draw the 4th shape.}
    \ealpara{(b) Ile zapałek używa?}{(b)\;How many matchsticks does it use?}
    \ealpara{(c) Którą tabliczkę mnożenia reprezentuje ten \ealkeytr{ciąg}?}{(c)\;Which \ealkey{times tables} does the \ealkey{sequence} represent?}
    \ashow{\ealpara{(b) 12 zapałek}{(b)\;12 matchsticks}}
    \ashow{\ealpara{(c) Tabliczkę mnożenia przez 3}{(c)\;The 3 \ealkey{times table}}}
  \end{minipage}\hfill
  \begin{minipage}[t]{0.48\textwidth}
    \small
    \ealpara{Q2. Uzupełnij \ealkeytr{ciąg}}{\textbf{Q2.\enspace Complete the \ealkey{Sequence}}}
    \footnotesize
    4,\;8,\;12,\;16,\;
    \ablank{20},\;\ablank{24},\;\ablank{28}
  \end{minipage}
\end{frame}

\begin{frame}[shrink=5]{Starter 6 (2 of 3)}
  \small
  \ealpara{Q3. Tabliczka mnożenia przez 3 ma \ealkeytr{wzór na n-ty wyraz} $3n$, jaki jest \ealkeytr{wzór na n-ty wyraz} dla tabliczki mnożenia przez 5?}{\textbf{Q3.\enspace The 3 \ealkey{times table} has \ealkey{$n$th term} rule $3n$, what is the \ealkey{$n$th term} rule for the 5 \ealkey{times table}?}}
  \footnotesize
  \ablank{$5n$}

  \vspace{0.45cm}

  \small
  \ealpara{Q4. Zapisz pierwsze 5 \ealkeytr{wyrazów} $3n+4$}{\textbf{Q4.\enspace Write the first 5 \ealkey{terms} of $\boldsymbol{3n+4}$}}
  \footnotesize
  \ablank{7},\;\ablank{10},\;\ablank{13},\;\ablank{16},\;\ablank{19}
\end{frame}

\begin{frame}[shrink=5]{Starter 6 (3 of 3)}
  \small
  \ealpara{Q5.}{\textbf{Q5.\enspace}}
  \ealpara{\ealkeytr{Ciąg} jest tworzony przez wzięcie \textbf{tabliczki mnożenia przez 3} i \textbf{odjęcie 2 od każdego \ealkeytr{wyrazu}}.}{A \ealkey{sequence} is made by taking the \textbf{3 \ealkey{times table}} and \textbf{subtracting 2 from every \ealkey{term}.}}
  \footnotesize
  \ealpara{(a) Zapisz pierwsze 5 \ealkeytr{wyrazów}.}{(a)\;Write the first 5 \ealkey{terms}.}
  \ealpara{(b) Jaki jest \ealkeytr{wzór na n-ty wyraz}?}{(b)\;What is the \ealkey{$n$th term} rule?}
  \ealpara{(c) Użyj swojej reguły, aby znaleźć 20. \ealkeytr{wyraz}.}{(c)\;Use your rule to find the \textbf{20th} \ealkey{term}.}
  \ealpara{(d) Czy $\boldsymbol{100}$ jest w \ealkeytr{ciągu}? Pokaż wyraźnie swoje obliczenia.}{(d)\;Is $\boldsymbol{100}$ in the \ealkey{sequence}? Show your working clearly.}
  \ashow{\ealpara{(a) $1, 4, 7, 10, 13$. \quad (b) $\boldsymbol{3n-2}$. \quad (c) $3\times20-2 = 58$.}{(a)\;1,\,4,\,7,\,10,\,13. \quad (b)\;$\boldsymbol{3n-2}$. \quad (c)\;$3\times20-2 = 58$.}}
  \ashow{\ealpara{(d) $3n-2=100 \Rightarrow 3n=102 \Rightarrow n=34$. \textbf{Tak}, 100 to 34. \ealkeytr{wyraz}.}{(d)\;$3n-2=100 \Rightarrow 3n=102 \Rightarrow n=34$. \textbf{Yes}, 100 is the 34th \ealkey{term}.}}
\end{frame}

% ================================================================
% STARTER 7
% ================================================================
\begin{frame}[shrink=5]{Starter 7 (1 of 3)}
  \begin{minipage}[t]{0.48\textwidth}
    \small
    \ealpara{Q1. Kropki}{\textbf{Q1.\enspace Dots}}
    \fourtabledots[red]{1}\quad
    \fourtabledots[blue]{2}\quad
    \fourtabledots[green!70!black]{3}\quad
    \ashowq{\fourtabledots[orange]{4}}\\[0.4em]
    \footnotesize
    \ealpara{(a) Narysuj 4. obrazek.}{(a)\;Draw the 4th picture.}
    \ealpara{(b) Ile kropek zawiera?}{(b)\;How many dots does it contain?}
    \ealpara{(c) Którą tabliczkę mnożenia reprezentuje ten \ealkeytr{ciąg}?}{(c)\;Which \ealkey{times table} does the \ealkey{sequence} represent?}
    \ashow{\ealpara{(b) 16 kropek}{(b)\;16 dots}}
    \ashow{\ealpara{(c) Tabliczkę mnożenia przez 4}{(c)\;The 4 \ealkey{times table}}}
  \end{minipage}\hfill
  \begin{minipage}[t]{0.48\textwidth}
    \small
    \ealpara{Q2. Uzupełnij \ealkeytr{ciąg}}{\textbf{Q2.\enspace Complete the \ealkey{Sequence}}}
    \footnotesize
    5,\;10,\;15,\;20,\;
    \ablank{25},\;\ablank{30},\;\ablank{35}
    \ashow{\ealpara{(dodaj 5 za każdym razem --- tabliczka mnożenia przez 5)}{(add 5 each time --- the 5 \ealkey{times table})}}
  \end{minipage}
\end{frame}

\begin{frame}[shrink=5]{Starter 7 (2 of 3)}
  \small
  \ealpara{Q3. Zapisz pierwsze 5 \ealkeytr{wyrazów} $4n-3$}{\textbf{Q3.\enspace Write the first 5 \ealkey{terms} of $\boldsymbol{4n-3}$}}
  \footnotesize
  \ablank{1},\;\ablank{5},\;\ablank{9},\;\ablank{13},\;\ablank{17}
\end{frame}

\begin{frame}[shrink=5]{Starter 7 (3 of 3)}
  \small
  \ealpara{Q4.}{\textbf{Q4.\enspace}}
  \ealpara{Zacznij od \textbf{tabliczki mnożenia przez 5} i \textbf{dodaj 1} do każdego \ealkeytr{wyrazu}.}{Start with the \textbf{5 \ealkey{times table}} and \textbf{add 1} to every \ealkey{term}.}
  \footnotesize
  \ealpara{(a) Zapisz pierwsze 5 \ealkeytr{wyrazów}.}{(a)\;Write the first 5 \ealkey{terms}.}
  \ealpara{(b) Jaki jest \ealkeytr{wzór na n-ty wyraz}?}{(b)\;What is the \ealkey{$n$th term} rule?}
  \ealpara{(c) Znajdź 20. \ealkeytr{wyraz}.}{(c)\;Find the 20th \ealkey{term}.}
  \ealpara{(d) Czy $\boldsymbol{104}$ jest w \ealkeytr{ciągu}? Wyjaśnij, skąd to wiesz.}{(d)\;Is $\boldsymbol{104}$ in the \ealkey{sequence}? Explain how you know.}
  \ashow{\ealpara{(a) $6, 11, 16, 21, 26$. \quad (b) $\boldsymbol{5n+1}$. \quad (c) $5\times20+1 = 101$.}{(a)\;6,\,11,\,16,\,21,\,26. \quad (b)\;$\boldsymbol{5n+1}$. \quad (c)\;$5\times20+1 = 101$.}}
  \ashow{\ealpara{(d) $5n+1=104 \Rightarrow 5n=103 \Rightarrow n=20.6$. \textbf{Nie}, 104 nie jest \ealkeytr{wyrazem}.}{(d)\;$5n+1=104 \Rightarrow 5n=103 \Rightarrow n=20.6$. \textbf{No}, 104 is not a \ealkey{term}.}}
\end{frame}

\end{document}